\subsubsection*{Solution}
\paragraph*{Step-by-Step Derivation}

We interpret "only two adjoint fields" literally: the independent fermionic letters are

\[
\psi,\qquad Q(\psi)=1,
\]

and

\[
\partial\psi,\qquad Q(\partial\psi)=2,
\]

with no fields $\partial^k\psi$ for $k\geq2$.

\subparagraph*{1. Single-letter index}

Both letters are fermionic, so each enters the Witten index with a minus sign:

\[
i(q)=-q-q^2.
\]

For adjoint letters, the gauge-singlet index is given by the Molien-Weyl integral \href{\detokenize{https://arxiv.org/html/2605.08085}}{as in Eq. (2) of this paper}:

\[
I_N(q)=\int_{U(N)}dU\, \exp\left[ \sum_{m=1}^{\infty}\frac{i(q^m)}{m} \operatorname{Tr}U^m\operatorname{Tr}U^{-m} \right].
\]

\subparagraph*{2. Definition of the trace-relation index}

Write

\[
I_N(q)=B_N(q)-F_N(q),
\]

where $B_N$ and $F_N$ count bosonic and fermionic gauge invariants. The numbers of finite-rank relations are

\[
r_B=B_\infty-B_N,\qquad r_F=F_\infty-F_N.
\]

Therefore their index is

\[
\mathcal R_N(q)=r_B-r_F =I_\infty(q)-I_N(q).
\]

This agrees with the convention that trace relations are operators present at infinite rank but removed at finite rank, used in the finite-rank literature \href{\detokenize{https://arxiv.org/html/2605.08085}}{here}.

\subparagraph*{3. Infinite-rank index}

At infinite rank, the multi-trace formula is \href{\detokenize{https://arxiv.org/html/2605.08085}}{Eq. (34) here}

\[
I_\infty(q)=\prod_{n=1}^{\infty}\frac{1}{1-i(q^n)}.
\]

Using $i(q^n)=-q^n-q^{2n}$,

\[
I_\infty(q) = \prod_{n=1}^{\infty} \frac{1}{1+q^n+q^{2n}}.
\]

Since

\[
1+x+x^2=\frac{1-x^3}{1-x},
\]

we obtain

\[
\begin{aligned}
I_\infty(q)
=
\prod_{n=1}^{\infty}
\frac{1-q^n}{1-q^{3n}}
=
\prod_{\substack{n\geq1\\3\nmid n}}(1-q^n).
\end{aligned}
\]

Up to charge $15$, only the factors

\[
n=1,2,4,5,7,8,10,11,13,14
\]

contribute. Multiplying them gives

\[
\begin{aligned}
I_\infty(q)
={}&1-q-q^2+q^3-q^4+2q^6-q^7-q^8+3q^9\\
&-2q^{10}-q^{11}+4q^{12}-3q^{13}-2q^{14}
+5q^{15}+O(q^{16}).
\end{aligned}
\]

\subparagraph*{4. Exact $U(2)$ index}

Let $z=x_1/x_2$ be the ratio of the two $U(2)$ eigenvalues. The adjoint weights are

\[
1,\quad 1,\quad z,\quad z^{-1}.
\]

For a fermionic letter of charge $r$, exponentiating its contribution gives

\[
\exp\left[ -\sum_{m=1}^{\infty} \frac{q^{rm}}{m}\left(\frac{x_i}{x_j}\right)^m \right] = 1-q^r\frac{x_i}{x_j}.
\]

Thus the $U(2)$ Weyl integral becomes the constant term

\[
\begin{aligned}
I_2(q)
={}&\frac12\operatorname{CT}_{z}\Big[
(1-z)(1-z^{-1})(1-q)^2(1-q^2)^2\\
&\hspace{29mm}\times
(1-qz)(1-qz^{-1})
(1-q^2z)(1-q^2z^{-1})
\Big].
\end{aligned}
\]

Define

\[
A(z)=(1-qz)(1-q^2z) =1-(q+q^2)z+q^3z^2.
\]

With

\[
s=q+q^2,\qquad p=q^3,
\]

the relevant Laurent coefficients of $A(z)A(z^{-1})$ are

\[
[z^0]A(z)A(z^{-1})=1+s^2+p^2,
\]

and

\[
[z^1]A(z)A(z^{-1})=-s(1+p).
\]

Since

\[
(1-z)(1-z^{-1})=2-z-z^{-1},
\]

we find

\[
\frac12\operatorname{CT}_{z} \left[(2-z-z^{-1})A(z)A(z^{-1})\right] = 1+s^2+p^2+s(1+p).
\]

Consequently,

\[
\begin{aligned}
I_2(q)
={}&(1-q)^2(1-q^2)^2
\left[1+(q+q^2)^2+q^6+(q+q^2)(1+q^3)\right]\\
={}&1-q-q^2+q^3-q^4+2q^6-q^8+q^9\\
&-q^{10}-q^{11}+q^{12}.
\end{aligned}
\]

There are no terms above $q^{12}$ because each adjoint fermion has $2^2=4$ Grassmann components, so the maximum possible charge is

\[
4(1)+4(2)=12.
\]

\subparagraph*{5. Subtraction}

Finally,

\[
\mathcal R_2(q)=I_\infty(q)-I_2(q).
\]

Term-by-term subtraction gives

\[
\mathcal R_2(q) = -q^7+2q^9-q^{10}+3q^{12} -3q^{13}-2q^{14}+5q^{15} +O(q^{16}).
\]

A negative coefficient means a net excess of fermionic trace relations, while a positive coefficient means a net excess of bosonic trace relations.
